%; whizzy paragraph -pdf xpdf -latex ./whizzypdfptex.sh
%; whizzy-paragraph "^\\\\begin{frame}"
% latex beamer presentation.
% platex, latex-beamer でコンパイルすることを想定。 

%     Tokyo Debian Meeting resources
%     Copyright (C) 2009 Junichi Uekawa
%     Copyright (C) 2009 Nobuhiro Iwamatsu

%     This program is free software; you can redistribute it and/or modify
%     it under the terms of the GNU General Public License as published by
%     the Free Software Foundation; either version 2 of the License, or
%     (at your option) any later version.

%     This program is distributed in the hope that it will be useful,
%     but WITHOUT ANY WARRANTY; without even the implied warreanty of
%     MERCHANTABILITY or FITNESS FOR A PARTICULAR PURPOSE.  See the
%     GNU General Public License for more details.

%     You should have received a copy of the GNU General Public License
%     along with this program; if not, write to the Free Software
%     Foundation, Inc., 51 Franklin St, Fifth Floor, Boston, MA  02110-1301 USA

\documentclass[cjk,dvipdfmx,12pt]{beamer}
\usetheme{Tokyo}
\usepackage{monthlypresentation}

%  preview (shell-command (concat "evince " (replace-regexp-in-string "tex$" "pdf"(buffer-file-name)) "&")) 
%  presentation (shell-command (concat "xpdf -fullscreen " (replace-regexp-in-string "tex$" "pdf"(buffer-file-name)) "&"))
%  presentation (shell-command (concat "evince " (replace-regexp-in-string "tex$" "pdf"(buffer-file-name)) "&"))

%http://www.naney.org/diki/dk/hyperref.html
%日本語EUC系環境の時
\AtBeginDvi{\special{pdf:tounicode EUC-UCS2}}
%シフトJIS系環境の時
%\AtBeginDvi{\special{pdf:tounicode 90ms-RKSJ-UCS2}}

\title{ldapvi \& python-ldap で stress-free life}
%\subtitle{第98回 2013年3月度}
\author{まえだこうへい mkouhei@debian.or.jp\\IRC nick: mkouhei}
\date{2013年3月16日}
\logo{\includegraphics[width=8cm]{image200607/openlogo-light.eps}}

\begin{document}

\frame{\titlepage{}}


\begin{frame}
 \frametitle{ナニそれおいしいの？}
  \begin{itemize}
   \item 知らずに恩恵を受けているのですよ
   \item JPのアカウント管理もLDAP使っているのですよ
  \end{itemize}
\end{frame}

\begin{frame}
{\Huge LDAP管理者100人に聞きました\footnote{脳内のべ人数}}
\end{frame}

\begin{frame}{LDAP管理の何がストレスですか？}
 \begin{enumerate}
  \item ？
  \item ？
  \item ？
 \end{enumerate}
\end{frame}


\begin{frame}{LDAP管理の何がストレスですか？}
 \begin{enumerate}
  \item LDIFまじやべぇ
  \item 設定方法がDBに変わっちゃうとか意味分からないんですけど？！
  \item どこがLight weightなんや\footnote{X.500に対してというマジレス不要}
 \end{enumerate}
\end{frame}

\begin{frame}
{\Huge 実際のとこ、どうですか？}
\end{frame}

\begin{frame}{ストレスを取り除くたった三つの方法\footnote{言いたかっただけ}}
 \begin{itemize}
  \item LDAPを止める
  \item LDIFを自分でパースしない、作らない
  \item LDAPに代わる仕組みを作る
 \end{itemize}
\end{frame}

\begin{frame}
{\Huge 今日は二つ目の話}
\end{frame}

\begin{frame}{LDIFをパースしない、作らないようにするには}
 \begin{itemize}
  \item ldapviをつかう
  \item LDAP bindingを使う
 \end{itemize}
\end{frame}

\begin{frame}{ldapviとは}
  \begin{itemize}
    \item viインタフェースでLDAPを操作するクライアントツール
  \end{itemize}
\end{frame}

\begin{frame}{ldapviのメリット}
  \begin{itemize}
    \item LDIFを自分で作る必要がない
    \item visudo(8)のように、変更前にschemaチェックされる
    \item 変更内容からLDIFを生成できる
  \end{itemize}
\end{frame}

\begin{frame}{LDAP bindingのあるプログラミング言語\footnote{Debianパッケージにあるもの}}
  \begin{itemize}
    \item Java, PHP, Ruby, python, lua, Erlang, Haskell, OCaml, Perl, ...\footnote{順不同}
    \item とりあえず大抵あるので好きなの使えばよいでしょう
  \end{itemize}
\end{frame}


\begin{frame}

{\Huge 今日はpython}

\end{frame}

\begin{frame}[containsverbatim]{python-ldapで扱うデータ構造}
  \begin{commandline}
[('DN', {'attribute': ['value'],
          'attribute': ['value', 'value', ...], ...}),
 ('DN', {'attribute': ['value'],
          'attribute': ['value', 'value', ...], ...}),
 ...
]
  \end{commandline}
\end{frame}

\begin{frame}

{\Huge ドキュメントが整っているので読めばだいたいわかる\footnote{一部間違っているけど}}

\end{frame}

\begin{frame}{さらに楽にするには}
\begin{itemize}
\item ldapvi \& debconf で LDAPサーバの構築自動化
\item プラス python-ldap \& CIで自分で開発しているLDAPソフトウェアをクリーン環境でテストを実行
\end{itemize}
\end{frame}

\begin{frame}

{\Huge 何か質問ある？}

\end{frame}

\end{document}
